% !PRINTED 628
\weberpage{628}
\begin{center}
\S.\,173.

\gesperrt{Representation of the prime factors of $p$.}
\end{center}

For the representation of the prime functionals of the field $\Omega_m$ we can apply a procedure that we became acquainted with in \S.\,157, and which, in consequence of the circumstance that
\[ 1, r, r^2, \ldots, r^{\mu-1} \]
is a basis of $\mathfrak{o}$, is here essentially simplified.

The natural prime $q$ is, as we have already seen, the $m$th power of a prime $\sigma$ existing in the field $\Omega_m$; with this we need not concern ourselves further, and we therefore consider here only the primes $p$ different from $q$.
% [sic] Source prints the m-th power of sigma here; the ramification index of the prime above q is mu = phi(m). Preserve the printed m.
Let $\mathfrak{p}$ be a prime factor of $p$, and let $\mathfrak{A}$ have the same meaning as above. If $p$ belongs to the exponent $f$ and
\[ ef = \varphi(m) = \mu \]
holds, then the group $\mathfrak{A}$ consists of the numbers
\[ 1, p, p^2, \ldots, p^{f-1}. \]
If now
\[ f(t) = N(t-r) = \prod^n (t - r^n) \]
denotes the polynomial in $t$ of degree $\mu$ whose roots are the $\mu$ quantities $r^n$, then we can decompose it into factors in such a way that each factor corresponds to one of the subsidiary groups of $\mathfrak{A}$. Namely, if we set
\weq{1}{F_1(t) = \prod_{0,\,f-1}^{h} (t - r^{p^h}),}
and in general, for every $n$ not divisible by $q$,
\weq{2}{F_n(t) = \prod_{0,\,f-1}^{h} (t - r^{np^h}),}
then $F_n(t)$ is identical with $F_{n'}(t)$ when $n$ and $n'$ belong to the same subsidiary group $\mathfrak{A}\xi_i$. If, then, we understand by $\xi_1, \xi_2, \ldots, \xi_e$ the system of numbers applied in \S.\,172, (5), we have
\weq{3}{f(t) = F_{\xi_1}(t)\, F_{\xi_2}(t) \ldots F_{\xi_e}(t).}

Now, however, by the binomial theorem:
\[ (t - r^{p^h})^p \equiv (t^p - r^{p^{h+1}}) \, (\mathrm{mod}\ p), \]
%CONT>>
